\section{Methodology}

\subsection{Theoretical Foundation}

We ground SSI in learning theory: models trained on diverse phrasings of the same content develop representations invariant to surface form. This invariance should manifest as uniform confidence scores across paraphrases. The causal chain is:

\begin{enumerate}
    \item Contamination exposure $\rightarrow$ Training includes benchmark items
    \item Representation invariance $\rightarrow$ Diverse exposure creates robust encoding
    \item Confidence uniformity $\rightarrow$ Invariant representations yield consistent predictions
    \item SSI signal $\rightarrow$ Low variance indicates potential contamination
\end{enumerate}

\subsection{SSI Formulation}

For each benchmark item, we generate $K$ paraphrases and extract model confidence scores:
\begin{equation}
\text{SSI}(x) = \frac{1}{\text{Var}(\{c(x_1), c(x_2), \ldots, c(x_K)\})}
\end{equation}

\subsection{Controlled Contamination}

We create contaminated model variants via LoRA fine-tuning of Mistral-7B-v0.1 at levels 0\%, 5\%, 10\%, 20\%, 50\% of MMLU test items.

\subsection{Paraphrase Generation}

$K=20$ paraphrases per item via T5-paraphrase, GPT-4, and rule-based methods.

\subsection{Hypothesis Decomposition}

\begin{table}[h]
\centering
\begin{tabular}{@{}llll@{}}
\toprule
ID & Type & Statement & Gate \\
\midrule
H-E1 & Existence & SSI differs between clean and contaminated models & MUST\_WORK \\
H-M1 & Mechanism & Contamination injection creates training exposure & MUST\_WORK \\
H-M2 & Mechanism & Training develops representation invariance & SHOULD\_WORK \\
H-M3 & Mechanism & Invariance manifests as confidence uniformity & SHOULD\_WORK \\
H-M4 & Mechanism & SSI captures invariance as contamination signal & SHOULD\_WORK \\
\bottomrule
\end{tabular}
\end{table}
